The implied volatility surface is the central object of equity derivatives markets, yet it is only
ever observed as a sparse, noisy, irregular cloud of option quotes. Reconstructing a complete,
smooth, arbitrage-free surface from these observations is an ill-posed inverse problem whose
solution feeds pricing, hedging and risk management. This thesis studies the problem across the
full methodological spectrum on a single unified S\&P~500 dataset (OptionMetrics, 2018--2025):
the parametric benchmark of Gatheral--Jacquier (SVI per slice, SSVI surface-wide, with their
complete no-arbitrage theory replicated and validated against the original paper), a
derivative-constrained deep smoother in the spirit of Ackerer et al.\ and Hoshisashi et al.\
(an SSVI prior multiplied by a neural corrector, no-arbitrage enforced through exact automatic
differentiation), and neural operators following Operator Deep Smoothing, which learn the
smoothing map itself and apply to unseen days without retraining. All methods share one
evaluation protocol: hold-out reconstruction error, a static-arbitrage audit built on the
Durrleman condition, robustness to quote sparsity, and behaviour across market regimes.
The thesis's central contribution is methodological and transversal: \emph{soft no-arbitrage
penalties are enforced only where they are sampled}. We exhibit a collocation blind spot in
which a calendar penalty is \emph{exactly} zero at its training nodes while an independent audit
grid reveals a $28\%$ violation rate between them; we trace it to the interaction between
exponential maturity spacing and the geometry of calendar arbitrage; we repair it with
constraint-specific collocation grids, reducing material violations from $34.4\%$ to $0.0\%$ on a
controlled stress test whose validity is itself verified by a two-condition gate; and we then show
the same pathology---in a stronger, two-dimensional form, blind between nodes \emph{and} across
days---in the neural operator, where the penalty-grid audit reports near-zero violations while
the independent audit flags $330$ of $386$ out-of-sample days as materially butterfly-violating.
The Gatheral--Jacquier conditions are used constructively: our SSVI prior is \emph{certified}
statically arbitrage-free on its working domain rather than luckily so. Finally, the audit is
priced: through the identity $\sigma_{\mathrm{loc}}^2 = \partial_\tau w / g$, arbitrage violations
become negative local variances (up to $50\%$ of the operator's surface requires repair before a
single Monte-Carlo path can be drawn) and negative risk-neutral density mass (up to $230\%$ on
some slices), and a residual-economics pipeline shows that surface dirtiness measurably degrades
the information content of quote--model residuals. The conclusion is a change of question: not
\emph{which model fits best}, but \emph{what a model's constraints actually guarantee, on which
set, and at what economic cost when they silently fail}.
